%=========== ΛΥΣΗ - ΘΕΜΑ Γ ===========

\begin{thema}{Γ}

\begin{center}
\begin{tikzpicture}[scale=0.75]
  \draw[help lines, gray!30] (-1,-1) grid (6,5);
  \draw[-latex] (-1.3,0) -- (6.3,0) node[right] {$x$};
  \draw[-latex] (0,-1.3) -- (0,5.3) node[above] {$y$};
  \foreach \x in {1,2,3,4,5} \draw (\x,0.1) -- (\x,-0.1) node[below,font=\footnotesize] {$\x$};
  \foreach \y in {1,2,3,4} \draw (0.1,\y) -- (-0.1,\y) node[left,font=\footnotesize] {$\y$};
  
  \draw[dashed, gray!50] (-1,-1) -- (5,5) node[right, font=\footnotesize, fill=white, inner sep=1pt] {$y=x$};
  
  \draw[very thick, blue!70!black, smooth, tension=0.4] plot coordinates {(0,0) (1,1) (2,1.6) (3,2) (4,3) (5,4)};
  \node[blue!70!black, fill=white, inner sep=1pt] at (4.5,2) {$C_f \ \text{(Κοίλη)}$};
  
  \draw[very thick, orange!80!black, smooth, tension=0.4] plot coordinates {(0,0) (1,1) (1.6,2) (2,3) (3,4) (4,5)};
  \node[orange!80!black, fill=white, inner sep=1pt] at (1,3.5) {$C_{f^{-1}} \ \text{(Κυρτή)}$};
  
  \filldraw[blue!70!black] (0,0) circle (2pt);
  \filldraw[blue!70!black] (1,1) circle (2pt) node[below right,font=\footnotesize] {$(1,1)$};
  \filldraw[blue!70!black] (3,2) circle (2pt) node[below,font=\footnotesize] {$(3,2)$};
  \filldraw[blue!70!black] (5,4) circle (2pt) node[below,font=\footnotesize] {$(5,4)$};
  
  \filldraw[orange!80!black] (2,3) circle (2pt) node[above,font=\footnotesize] {$(2,3)$};
  \filldraw[orange!80!black] (4,5) circle (2pt) node[above,font=\footnotesize] {$(4,5)$};
  
  \fill[yellow!15, opacity=0.5, smooth, tension=0.4]
    plot coordinates {(1,1) (2,1.6) (3,2) (4,3)}
    -- plot[smooth, tension=0.4] coordinates {(4,3) (3,4) (2,3) (1.6,2) (1,1)} -- cycle;

  % Arrows indicating symmetry
  \draw[latex-latex, dashed, purple!70!black] (3,2) -- (2,3);
  \draw[latex-latex, dashed, purple!70!black] (5,4) -- (4,5);
  
  \node[purple!70!black, font=\footnotesize, fill=white, inner sep=1pt] at (2.4, 4.3) {$\text{Συμμετρία ως προς } y=x$};
\end{tikzpicture}
\end{center}

\begin{erwthma}
    \item Από τα δεδομένα της εκφώνησης και τη γραφική παράσταση της $C_f$, έχουμε:
    \[ f(3) = 2 \implies f^{-1}(2) = 3 \]
    \[ f(5) = 4 \implies f^{-1}(4) = 5 \]
    Επιπλέον, εφόσον η $f^{-1}$ ορίζεται στο διάστημα $[0, 4]$ και το $3 \in [0, 4]$, ισχύει από τον ορισμό της αντίστροφης συνάρτησης ότι:
    \[ (f \circ f^{-1})(3) = 3 \]

    \item Γνωρίζουμε ότι για μια γνησίως αύξουσα συνάρτηση $f:[a, b] \to [f(a), f(b)]$, ισχύει η σχέση εμβαδών:
    \[ \int_{a}^{b} f(x) \, \d x + \int_{f(a)}^{f(b)} f^{-1}(y) \, \d y = b \cdot f(b) - a \cdot f(a) \]
    Εφαρμόζοντας τη σχέση για $a=0$ και $b=5$, με $f(0)=0$ και $f(5)=4$, έχουμε:
    \[ \int_{0}^{5} f(x) \, \d x + \int_{0}^{4} f^{-1}(y) \, \d y = 5 \cdot 4 - 0 \cdot 0 = 20 \]
    Δεδομένου ότι $\int_{0}^{5} f(x) \, \d x = 9$, έπεται:
    \[ 9 + \int_{0}^{4} f^{-1}(y) \, \d y = 20 \implies \int_{0}^{4} f^{-1}(y) \, \d y = 11 \]

    \item \textbf{Γραφικά:} Η γραφική παράσταση $C_{f^{-1}}$ προκύπτει από την συμμετρία της $C_f$ ως προς την ευθεία $y=x$. Εφόσον η $f$ είναι \textbf{κοίλη} στο $[0, 5]$, η $C_f$ «καμπυρίζει» προς τα κάτω. Λόγω της συμμετρίας, η $C_{f^{-1}}$ θα «καμπυρίζει» προς τα πάνω, άρα η $f^{-1}$ είναι \textbf{κυρτή}.
    
    \textbf{Αλγεβρικά:} Έστω $g(y) = f^{-1}(y)$. Η παράγωγος της $g$ είναι $g'(y) = \frac{1}{f'(g(y))}$. Παραγωγίζοντας ξανά ως προς $y$:
    \[ g''(y) = \frac{d}{dy} [f'(g(y))]^{-1} = -[f'(g(y))]^{-2} \cdot f''(g(y)) \cdot g'(y) = -\frac{f''(g(y)) \cdot g'(y)}{[f'(g(y))]^2} \]
    Επειδή η $f$ είναι γνησίως αύξουσα, ισχύει $f'(x) > 0 \implies g'(y) > 0$. Επειδή η $f$ είναι κοίλη, ισχύει $f''(x) \le 0$. Συνεπώς:
    \[ g''(y) = -\frac{(\le 0)(+)}{(+)} \ge 0 \]
    Άρα η $f^{-1}$ είναι \textbf{κυρτή}.

    \item Έστω $x(t)$ και $y(t)$ οι συντεταγμένες του σημείου στην $C_{f^{-1}}$. Γνωρίζουμε ότι $x = f^{-1}(y)$.
    Ο ρυθμός μεταβολής της τετμημένης είναι:
    \[ \frac{\d x}{\d t} = \frac{\d}{\d t} f^{-1}(y) = (f^{-1})'(y) \cdot \frac{\d y}{\d t} \]
    Στη στιγμή που $y_0 = 3$, έχουμε από την $C_{f^{-1}}$ (ή από την $C_f$ αφού $f(4)=3$) ότι $x_0 = f^{-1}(3) = 4$. 
    Χρησιμοποιώντας τη σχέση $(f^{-1})' = \frac{1}{f'}$, έχουμε:
    \[ \frac{\d x}{\d t} = \frac{1}{f'(x_0)} \cdot \frac{\d y}{\d t} = \frac{1}{f'(4)} \cdot 3 \]
    Από τα σημεία $(3,2), (4,3), (5,4)$ της $C_f$, παρατηρούμε ότι αυτά είναι συνευθειακά και ανήκουν στην ευθεία $y = x-1$. Έτσι, στο διάστημα $(3, 5)$, η $f'(x) = 1$. Άρα $f'(4) = 1$.
    \[ \frac{\d x}{\d t} = \frac{1}{1} \cdot 3 = 3 \text{ μον/sec} \]
\end{erwthma}
\end{thema}

%=========== ΤΕΛΟΣ ΛΥΣΗΣ - ΘΕΜΑ Γ ===========